\section{讲解笔记：【第6次作业笔记】 组合结构(1)}
\begin{analysisbox}
本节按题目、答案和讲解笔记顺序整理。对原图中的结构几何关系，正文保留 可辨识文字，并在图形页中保留清理后的题图，便于核对杆件、支座、荷载和尺寸。
\end{analysisbox}
\subsection*{第 1 页转写}
结力基础班第6次作业  组合结构

组合结构作业及答案

C

烟

 FN

lil

FN.ws450，2 印

DEliidhoFNGF.in

Fn 髺



髻

堡

兽

武客29 答

Fsa

厅二恶

FNA13

2厅

一普\_\_E

E

Ftii



n.in

EME

hn

iti.if

4 fi 下二管

IMA 0

后E 3a fxa

fxzaioi.FD.EE



Fm答

Minfix za in a f a 0

特别注意：梁式杆中 轴力，剪力和弯矩，不要遗漏。

\subsection*{第 2 页转写}
2 - 35

[r3=0

2412

fnixl + 4x2lxl=0

Ful= -2 $\alpha$

EFy=D

29l+ Fm+ Fu2= 0, F2=0

图2-165

29

2-36

试作图2-166所示结构的弯矩图、剪力图、轴力图。（哈尔滨工业大学2010）

a=g xoz + \textasciicircum{} $\times$o + x9

.. Fox- -40 kw

4o 20kN

20kN/m

7140

20

140

3m

FVDA

350

" FWPA - KN

4m

08

c81

图2-166

130

30kn

1 30

个50

130

30kn

350

\subsection*{第 3 页转写}
2 - 37

作图2-167所示结构的弯矩图。（重炭

1Fp

ZMA-0

a=17 * - 12 $\times$861

图2-167

2 - 38

作图2－168所示结构的弯矩图。（重庆大学2012）

2kN/m

Fex .8 -2xtx b:0

4m

" FEx= 6

4m

<FEX=6

m图(kw.)

4m

4m

图2-168

\subsection*{第 4 页转写}
2-39

FP

Fp!

.. MA= 2pL

FpL

忽略了BD中的剪力

MA

图2-169

\subsection*{第 5 页转写}
\emph{}

\subsection*{第 6 页转写}
\emph{}

